\noindent Grupo 29

\medskip
\noindent ¿Hicieron uso de herramientas de IA en su trabajo para este informe y/o presentación? \quad SI \underline{\ X\ } \quad NO \underline{\ \ \ }

\medskip
\noindent ¿Cuáles de las siguientes alternativas mejor describe el uso que hicieron de IA (puede ser más de una)?

\begin{itemize}
\item[X] Apoyo a la redacción, ortografía, estructura gramatical, etc., de los documentos.
\item[X] Apoyo al análisis de datos e información asociada al problema.
\item[X] Búsqueda y/o filtro de referencias, contenidos web, literatura, etc., relevante al proyecto.
\item[\ ] Apoyo para el desarrollo de modelos asociados al proyecto.
\item[X] Apoyo en el desarrollo de código computacional necesario para el proyecto.
\item[X] Apoyo en el análisis de resultados (por ejemplo, de distintas corridas computacionales, tendencias de predicciones, etc.).
\item[X] Apoyo para construir figuras, diagramas, etc.
\item[\ ] Apoyo para obtener conclusiones a partir de los análisis de resultados.
\item[\ ] Otro, especificar.
\end{itemize}

\noindent Herramientas usadas. Claude (Anthropic), a través de Claude Code.

\medskip
\noindent Reflexión. La IA fue un apoyo útil, sobre todo para escribir y depurar código, ordenar los resultados de muchas corridas y dar forma a las tablas y figuras. Después de las alucinaciones que detectamos en el informe~1, cada referencia se verificó en su fuente antes de citarla, y cada cifra del informe la genera un script desde los resultados, sin copiarla a mano. El planteamiento del problema, las decisiones de modelación, la elección del método y del caso base, y la interpretación de los resultados las tomó el grupo en sus reuniones y en las conversaciones con el profesor y el ayudante. Todo el código se revisó y se probó antes de usar sus resultados.

\vspace{1.2cm}
\noindent\begin{tabular}{@{}p{0.45\textwidth}p{0.45\textwidth}@{}}
\rule{5.5cm}{0.4pt} & \rule{5.5cm}{0.4pt} \\
Gonzalo Anguita & Diego Barrios \\[1cm]
\rule{5.5cm}{0.4pt} & \rule{5.5cm}{0.4pt} \\
Gustavo Cueto & Nicolás Farías \\[1cm]
\rule{5.5cm}{0.4pt} & \rule{5.5cm}{0.4pt} \\
Fernando González & Ricardo Palma \\[1cm]
\rule{5.5cm}{0.4pt} & \\
Cristóbal Rehbein & \\
\end{tabular}
